\chapter[Sermon 2. Of the Title of the Whole Work, and Exposition Thereof.]{Of the Title of the Whole Work, and Exposition Thereof.}
\label{sermon:002}
\markright{Sermon 2. Of the Title of the Whole Work, and Exposition Thereof.}

I said the whole book was contained in six parts. Three members of the first part. Now must we look on the first part: Which hath chiefly three members: The title, beginning, and brief narration. For this present we will only speak of the Title, which is thus.

The revelation of Jesus Christ, which God gave to him, to show his servants the things that must shortly come to pass. He sent and showed by his angel to his servant John, who bore witness to the word of God, and to the testimony of Jesus Christ, and to all things that he saw. Happy is he who reads, and those who hear the words of the prophecy, and keep those things written therein. For the time is at hand.

This title is abundant and proclaims all the beneficial matters that ought to be stated at the start of books.

\index{Apocalypse (Book of)}First, the title or inscription of the entire work is set: the Apocalypse, or revelation of Jesus Christ, which truly was opened or revealed by Jesus Christ himself.

This title immediately demonstrates that this work is not a human invention, but a divine doctrine: as that which was revealed by our Lord, king, and priest Jesus Christ, out of heaven, from the right hand of the Father, exercising there the office of the high priest, and still teaching us profitable things. And although it is also called the Revelation of John, it is ascribed to him for no other reason than that, as a scribe, he wrote and presented it.

Again, it is declared even more plainly from whence this Revelation comes: even from God himself. For he says, “which God,” namely the Father, gave unto him, that is, to Christ. For in the holy and blessed Trinity, there is a distinction of persons. And although all things of the Father are also the Son’s, and all things of the Son are likewise the Father’s, yet Scripture mentions the Father giving to the Son, and the Son receiving from the Father.

All the ancient writers have consistently explained this as occurring through the mystery of dispensation. For the Son received something from the Father, just as the Son of God himself says, “Glorify me with the glory that I had with you before the world existed.”

Moreover, the Son is the wisdom, word, and mouth of the Father, by whom God in times past and now spoke and speaks to the Fathers, Prophets, Apostles, and to the universal church. The Father, by dispensation, gave to his Son this office, that he should be Bishop. For no man has seen God at any time; the only-begotten Son, who is in the bosom of the Father, has revealed him to us. Let us know, therefore, that this same is a Divine Revelation, which God the Father, loving mankind, has revealed by the only Bishop Christ unto his Church. And so it joins together the Father and the Son, that nevertheless the holy distinction of persons remains safe.

Now also is added, for what purpose God the Father has revealed, or granted the gift of revealing, namely, the office of priesthood to his Son, our Lord Jesus Christ: so that, being revealed, he might demonstrate it, and as it were set it before the eyes of his servants, that is, his worshippers and Christians, who are called the servants of God for their willing obedience. And just as a servant owes to his lord all that he has or is worth, so we owe to God ourselves wholly, and all that belongs to us, or else we are free, and not bound.

It is also declared to whom this revelation is opened. To all the servants of God. If, therefore, you desire to be called a servant of God, hear this book and remember it. And know that this book is prepared for you by God.

After he summarizes in a few words what Christ revealed to John, the events that must shortly come to pass. Therefore, the destinies of the Church are recounted, what good and evil things will happen to the faithful, and likewise what punishments must be inflicted upon the wicked.

And let no one gather from this teaching a sense of obligation, necessity, as though God did not work freely; or that the wicked do evil, not through their own fault, but by God’s compulsion. Good things must be done, because God, willingly binding himself to us by his promise, cannot fail to do what he does and promises; nevertheless, he works freely.

\index{John, Saint}\index{Judgment, Last}Good things must be done in the godly community, for the nature of grace and faith is such, just as the characteristic of ungodliness is to despise and transgress. Therefore, they must also be punished. And because the world is as it is, there must needs be countless heresies and calamities. He says these things must shortly be accomplished, for certain things began in the very time of Saint John. And although many things are found to be done a thousand years after, Saint Peter, the apostle, says, “A thousand years before the Lord are as it were yesterday.” Therefore, this Apocalypse pertains to the times of the primitive and last church, and declares whatever things shall happen to it until the last judgment. Indeed, how it shall reign forever.

\index{Daniel (Book of)}\index{Hebrews, Epistle to the}Moreover, the manner of revealing is also addressed. For Christ revealed those things, sending by his angel, or his angel sent forth, to whom he gave commandment what he should say and do. Whereupon this angel is afterward also called Christ, because he represented the person of Christ. Therefore, we must not consider the angel in this book, but always Christ, the true author of all these things. Indeed, the divinity of Christ is commended to us when we hear that Christ is the Lord of angels. Paul reasoned more at length about this to the Hebrews. Moses in the twelfth chapter of Numbers sets forth chiefly three manners of prophecy or revelation. First, by vision, of which sort many are ascribed to Daniel, one notable to Peter in the tenth chapter of Acts, and likewise to Paul. And into this form falls also the Apocalypse. Secondly, by dream: of which sort were those of Pharaoh and Nebuchadnezzar, kings whom Joseph and Daniel interpreted. The Prophet Joel in the second chapter mentions visions and dreams. For in the New Testament also there are very many holy and prophetical dreams. Lastly, Moses sets forth a skillful exposition, as many were made to Moses and to the Apostles. Into whose fellowship the Apocalypse comes after a sort also, where visions are openly explained. Here appears the unspeakable goodness of God, which in many ways procures and works our salvation, and pleasantly prepared offers it unto us to enjoy. Unhappy is he who does not know these things.

Besides this, much mention is made of him to whom Christ has opened this divine and most excellent revelation, given to John. He commends himself (for it was expedient for confuting his adversaries, seeing that Paul also many times maintained his authority against the false Apostles) by four epithets. First, he calls himself the servant of Christ. This is the oldest and noblest title, which the fathers, prophets, and apostles have used, for they dedicate and consecrate themselves to God. Secondly, John testified to the word of God among the apostles, most expressly declaring the divinity of Christ, especially where he testified and said, “In the beginning was the word.” Moreover, he was a witness of the witness of Jesus Christ: under which name the Lord himself in the Gospel, and John in the twelfth chapter of his Gospel, comprised the whole evangelical doctrine. And he was a seeing witness of all these things. For in the first chapter we have seen his glory, he says, and in the nineteenth chapter he saw water and blood gush out of the Lord’s side. In his epistle, he says, “We have seen and have heard.”

\index{Arethas of Caesarea}Arethas notes that in some Greek manuscripts is added what is also found in the Greek manuscript of Spain. And whatever he has heard, and whatever may be, and whatever must be done afterward.

That same John is the author of this book, who, just as he saw the Lord in the flesh upon earth, so he saw him in spirit, revealing these things in heaven, and who presents to the church sights that are most certain and sure. This John was that beloved disciple of the Lord, who at the last supper rested upon his breast, to whom in his last will he bequeathed his mother on the cross, one virgin to another.

\index{Irenaeus}\index{Eusebius of Caesarea}He alone stood by at the altar of the cross when Christ died, a witness to the true death and to our purification. He lived until the time of Emperor Trajan, as Eusebius recounts in his chronicles, citing Irenaeus, noting the year as one hundred and three from the birth of Christ. Dorotheus, an ancient writer, affirms that John lived sixty years.

It is also important to consider the benefit of this godly work or revelation, so that readers and hearers might be spurred to diligence. For this book is also called a prophecy. This is because it foretells future events, making it the prophecy of the New Testament. Moreover, a prophecy is an exposition that clarifies and explains the Old Prophets. It promises blessedness to those who read, hear, and keep the things written in this book. Blessedness encompasses the benefits of the present life, to the extent that the Lord allows them to the faithful, but chiefly pertains to the life to come. If the benefit of this book has already been discussed in the first sermon. Note that it is not enough simply to read or hear this book; it must be put into practice and kept diligently.

\index{Antichrist}For the Lord also said in the Gospel: Blessed are those who hear the word of God and keep it. Therefore, those who shape their lives according to this book are blessed. For they avoid the deception of Antichrist, remain in the faith of Christ, and live forevermore. \&c.

And he concludes the title with an exclamation, by which he stirs up the listeners greatly: “The time is at hand!” As though he should say: Let no one think here that strange things, which concern him not, are told here, which shall come to pass at length after many worlds; they belong to every one of us. For they are written of matters that chiefly concern us, and even of our own affairs. Thus he shows that this book is profitable for all men, ages, and worlds. God the Father, by his Son, teaches profitable things, and admonishing in due season, be praised world without end. Amen.
